\documentclass{article}

\usepackage{amsmath,amssymb}
\usepackage{xurl}
\usepackage[hidelinks]{hyperref}
\setlength{\emergencystretch}{2em}

% Shared commands for notes in docs/paper-gaps/.

\DeclareUnicodeCharacter{2080}{\ifmmode{}_{0}\else\textsubscript{0}\fi}
\DeclareUnicodeCharacter{2081}{\ifmmode{}_{1}\else\textsubscript{1}\fi}
\DeclareUnicodeCharacter{2082}{\ifmmode{}_{2}\else\textsubscript{2}\fi}
\DeclareUnicodeCharacter{2099}{\ifmmode{}_{n}\else\textsubscript{n}\fi}

\newcommand{\Lean}{\textsc{Lean}}
\ExplSyntaxOn
\NewDocumentCommand{\leanid}{m}{
  \ifmmode
    \textsf{\detokenize{#1}}
  \else
    \group_begin:
    \urlstyle{sf}
    \tl_set:Nn \l_tmpa_tl {#1}
    \tl_replace_all:Nnn \l_tmpa_tl {\_} {_}
    \exp_args:NV \path \l_tmpa_tl
    \group_end:
  \fi
}
\ExplSyntaxOff
\providecommand{\tr}{\operatorname{tr}}
\newcommand{\tnleanIssueBase}{https://github.com/LionSR/TNLean/issues/}
\newcommand{\tnleanPrBase}{https://github.com/LionSR/TNLean/pull/}
\newcommand{\tnleanDocBase}{https://sirui-lu.com/TNLean/docs/}
\newcommand{\ghissue}[1]{\href{\tnleanIssueBase#1}{GitHub issue \##1}}
\newcommand{\ghpr}[1]{\href{\tnleanPrBase#1}{PR \##1}}
\newcommand{\doclink}[1]{\href{\tnleanDocBase#1.html}{\path{#1}}}

% Dirac notation and the identity operator, as in blueprint/src/macros/common.tex.
% \providecommand keeps note-local definitions authoritative.
\usepackage{dsfont}
\providecommand{\ket}[1]{|#1\rangle}
\providecommand{\bra}[1]{\langle #1|}
\providecommand{\braket}[2]{\langle #1 | #2 \rangle}
\providecommand{\ketbra}[2]{|#1\rangle\!\langle #2|}
\providecommand{\Id}{\mathds{1}}
\providecommand{\one}{\mathds{1}}
\providecommand{\C}{\mathbb{C}}
\providecommand{\MN}[1]{M_{#1}(\C)}
\providecommand{\MD}{M_D(\C)}

% External referencing for paper sources.  The build helper writes sanitized
% auxiliary files under build/paper-aux/ and excludes duplicated source labels.
\usepackage{xr}

% Import a generated paper auxiliary file with a caller-supplied label prefix.
% The candidate paths support builds from docs/paper-gaps/, the repository root,
% and build/paper-aux/ itself.
\newcommand{\importpaperaux}[2]{%
  \IfFileExists{../../build/paper-aux/#2.aux}{%
    \externaldocument[#1]{../../build/paper-aux/#2}%
  }{%
    \IfFileExists{build/paper-aux/#2.aux}{%
      \externaldocument[#1]{build/paper-aux/#2}%
    }{%
      \IfFileExists{#2.aux}{%
        \externaldocument[#1]{#2}%
      }{}%
    }%
  }%
}

% General external reference macros
\newcommand{\paperref}[2]{\ref{#1-#2}}
\newcommand{\papereqref}[2]{(\ref{#1-#2})}

% CPSV16 source (1606.00608 / MPDO-22-12-17-2.tex) external document and shortcuts
\importpaperaux{cpsv16-}{1606.00608/MPDO-22-12-17-2-tnlean}
\newcommand{\cpsvref}[1]{\paperref{cpsv16}{#1}}
\newcommand{\cpsveqref}[1]{\papereqref{cpsv16}{#1}}

% Verdict marker: \gapnote{<kind>}{<status>}. Kind is the policy
% classification (clarification, local-correction, scope-restriction,
% unfaithful, false-source, open-gap); status is open, wip (elimination
% underway), resolved, or historical. Machine-read by the site index;
% prints a verdict line.
\newcommand{\gapnote}[2]{\par\noindent\textbf{Verdict:} #1 (#2).\par}


\title{Diagonal Restriction and the Vertical Canonical Form in Proposition 4.13}
\author{}
\date{2026-07-11}

\begin{document}
\maketitle
\gapnote{unfaithful}{resolved}

\section{The source argument}

Proposition~4.13 of Cirac, P\'erez-Garc\'ia, Schuch, and Verstraete concerns a
tensor in horizontal canonical form whose matrix product operators are positive
semidefinite at every system size. It concludes that the same tensor, viewed
vertically, admits a vertical canonical form
\cite[Proposition~4.13]{Cirac2016MPDO_arXiv}. The local source is
\path{Papers/1606.00608/MPDO-22-12-17-2.tex}, lines~1863--1921.

The source does not restrict each horizontal normal tensor from physical pairs
$(i,j)$ to diagonal pairs $(i,i)$ and assert that the restriction remains
normal. Instead, it constructs a new canonical form in the vertical direction.
Positivity and the first-site insertion lemma exclude one-sided invariant
projections at lines~1873--1887 of
Ref.~\cite{Cirac2016MPDO_arXiv}. Positivity then excludes nontrivial periodic
sectors at lines~1889--1893. Only after these steps does the proof invoke a
canonical-form decomposition in the vertical direction. Lines~1895--1921 prove
that the resulting multiplicity weights are positive and that the associated
similarities reduce to unitary changes of basis.

The block-injectivity proposition at lines~340--345 of
Ref.~\cite{Cirac2016MPDO_arXiv} serves a different purpose: it separates the
normal blocks within one fixed canonical-form decomposition. It neither
compares horizontal blocks with their diagonal restrictions nor shows that
injectivity survives deletion of every off-diagonal physical matrix.

\section{A finite-dimensional obstruction}

The stronger diagonal-restriction assertion is false under the stated
horizontal canonical-form hypotheses.\footnote{The correction is tracked in
\href{https://github.com/LionSR/TNLean/issues/2537}{issue 2537}.}
Let $E_{ab}\in M_2(\mathbb C)$ denote the matrix units for
$a,b\in\{0,1\}$. Define
\begin{align}
    A^{(a,b)}
        &=c_aE_{ab},
    &
    c_0&=\frac{3}{5},
    &
    c_1&=\frac{4}{5},
    &
    a,b&\in\{0,1\}.
    \label{eq:cpgsv17_vertical_diagonal_restriction_matrix_units}
\end{align}
The four matrices in
Eq.~\ref{eq:cpgsv17_vertical_diagonal_restriction_matrix_units} are nonzero
scalar multiples of the matrix units and therefore span $M_2(\mathbb C)$.
They also satisfy
\begin{align}
    \sum_{a,b}
        (A^{(a,b)})^*A^{(a,b)}
        &=
        (c_0^2+c_1^2)I
        =I.
    \label{eq:cpgsv17_vertical_diagonal_restriction_left_canonical}
\end{align}
With one horizontal block, simultaneous block separation is ordinary
injectivity. Hence this example satisfies the injectivity, left-canonicality,
nonzero-weight, and finite-length block-separation assumptions of the older
per-copy \texttt{HorizontalCFData} structure.

Its diagonal restriction is
\begin{align}
    A^{\mathrm{diag}}(0)
        &=\frac{3}{5}E_{00},
    \label{eq:cpgsv17_vertical_diagonal_restriction_diag_zero}\\
    A^{\mathrm{diag}}(1)
        &=\frac{4}{5}E_{11}.
    \label{eq:cpgsv17_vertical_diagonal_restriction_diag_one}
\end{align}
Every word in the matrices from
Eqs.~\ref{eq:cpgsv17_vertical_diagonal_restriction_diag_zero}
and~\ref{eq:cpgsv17_vertical_diagonal_restriction_diag_one} is diagonal.
No blocking length can therefore span $M_2(\mathbb C)$, so the restricted
tensor is not normal. This example is formalized in
\path{TNLean/MPS/MPDO/BiCFDerivation/DiagonalRestrictionCounterexample.lean}.

This construction is not a counterexample to Proposition~4.13 of
Ref.~\cite{Cirac2016MPDO_arXiv}. The formal horizontal structure records
properties of a block decomposition but does not assume that its tensor
generates positive semidefinite operators at every length. That positivity
hypothesis is essential to the source proof. More importantly, even when
positivity holds, Proposition~4.13 constructs a new vertical basis of normal
tensors rather than reusing diagonal restrictions of the horizontal basis.

\section{The corrected formal theorem}

The literal theorem now assumes the CPSV canonical form of Proposition~4.13
together with the MPDO positivity hypothesis, and it concludes vertical
canonical form.\footnote{Formal declaration:
\leanid{MPOTensor.verticalCF\_of\_cpsvCanonicalForm} in
\doclink{TNLean/MPS/MPDO/CPSVVerticalCanonicalForm}.}
A separate theorem retains the normalized BNT-refined horizontal hypothesis
\path{MPOTensor.IsHorizontalCF} for arguments that already use this stronger
form.\footnote{Formal declaration:
\leanid{MPOTensor.verticalCF\_of\_horizontalCF} in
\doclink{TNLean/MPS/MPDO/VerticalCanonicalForm}.}

The proof follows the three stages of Proposition~4.13 in
Ref.~\cite{Cirac2016MPDO_arXiv}:
\begin{enumerate}
    \item positivity and the first-site insertion lemma exclude one-sided
    invariant projections in the vertical direction;
    \item positivity excludes nontrivial vertical periods;
    \item an independent vertical canonical-form decomposition supplies the
    vertical blocks, after which the proof establishes positivity of the
    multiplicity coefficients and unitary normalization of the sector gauges.
\end{enumerate}

The target is the vertically viewed tensor. Its letters are indexed by
horizontal bond pairs, and its matrices act on the physical space. In the
rectangular orientation used by the source, the displayed map is a coisometry
onto the retained sectors. The earlier statement about the diagonal tensor
$i\mapsto M^{ii}$ has been removed. The definitional realignment is documented
in \path{cpgsv17_vertical_tensor_definition.tex}.

The formal equality of matrix product vectors after diagonal restriction
remains useful as an auxiliary identity, but it cannot produce the normal
blocks of the vertical canonical form. Newton--Girard identities for scalar
weights concern only multiplicity coefficients; they cannot recover matrix
directions removed by $(i,j)\mapsto(i,i)$. The phase-class grouping of vertical
blocks and the resulting positive diagonal coefficients are described in
\path{cpgsv17_vertical_cf_grouping.tex}; they do not arise from diagonal
restriction.

The first operator identity from lines~1873--1887 of
Ref.~\cite{Cirac2016MPDO_arXiv} is formalized. If $H\geq0$, $P=P^*$, and
$PH=PHP$, then
\begin{align}
    HP
        &=(PH)^*
        =(PHP)^*
        =PHP,
    \label{eq:cpgsv17_vertical_diagonal_restriction_positive_corner_adjoint}\\
    (I-P)HP
        &=0.
    \label{eq:cpgsv17_vertical_diagonal_restriction_positive_corner_vanishes}
\end{align}
For doubled indices $p=(a,b)$ and $q=(c,e)$, define the three first-site
contractions by
\begin{align}
    (Y_L)_{p,q}
        &=P_{a,c}\delta_{b,e},
    \label{eq:cpgsv17_vertical_diagonal_restriction_left_insertion}\\
    (Y_R)_{p,q}
        &=\delta_{a,c}P_{e,b},
    \label{eq:cpgsv17_vertical_diagonal_restriction_right_insertion}\\
    (Y_C)_{p,q}
        &=P_{a,c}P_{e,b}.
    \label{eq:cpgsv17_vertical_diagonal_restriction_corner_insertion}
\end{align}
Their contractions with the horizontal matrix product vector are respectively
the entries of $PH^{(N)}$, $H^{(N)}P$, and $PH^{(N)}P$. The two applications
of the first-site insertion lemma are therefore formalized without abstract
placeholder matrices: the two entrywise operator identities force equality of
the opposite-corner insertions on every horizontal canonical-form block.

The first stage is complete. Let $P_1$ act as $P$ on the first site and as the
identity on all remaining sites:
\begin{align}
    P_1
        &=P\otimes I^{\otimes(N-1)}.
    \label{eq:cpgsv17_vertical_diagonal_restriction_first_site_projection}
\end{align}
The vertically viewed relation $PM=PMP$ gives
\begin{align}
    P_1H^{(N)}
        &=P_1H^{(N)}P_1
        =H^{(N)}P_1.
    \label{eq:cpgsv17_vertical_diagonal_restriction_first_site_reduction}
\end{align}
The representative-grouped Lemma~L transfers
Eq.~\ref{eq:cpgsv17_vertical_diagonal_restriction_first_site_reduction}
through every repeated weighted sector and through the direct sum of the copy
gauges. Thus the original tensor satisfies
\begin{align}
    MP
        &=PMP,
    \label{eq:cpgsv17_vertical_diagonal_restriction_opposite_corner_identity}\\
    (I-P)MP
        &=0.
    \label{eq:cpgsv17_vertical_diagonal_restriction_opposite_corner_vanishes}
\end{align}
The formal theorem assumes only $P=P^*$; the source specializes $P$ to an
orthogonal projector.

The final positive-operator implication needed for the periodic-sector stage
is also formalized:
\begin{align}
    Q[H^{(N)}]^p
        &=[H^{(N)}]^pQ
    \quad\Longrightarrow\quad
    QH^{(N)}=H^{(N)}Q.
    \label{eq:cpgsv17_vertical_diagonal_restriction_power_commutation}
\end{align}
A nontrivial vertical period produces a displaced cyclic projector and its word
invariance. The horizontal canonical-form version of Lemma~L supplies a chain
length at which this projector does not commute with the operator. Positivity
at that length and
Eq.~\ref{eq:cpgsv17_vertical_diagonal_restriction_power_commutation} yield the
contradiction. This excludes nontrivial vertical periods without the stronger
all-length inequality printed in Ref.~\cite{Cirac2016MPDO_arXiv}.

Invariant-projector closure is established for the vertically viewed tensor.
It produces a complete orthogonal family of irreducible reducing corners with
exact physical isometries in
\path{TNLean/MPS/MPDO/VerticalReduction.lean}. After exactly the zero corners
are discarded, the retained corners are spectrally normalized to normal
tensors. This construction preserves the isometries, both intertwining
identities, exact compression, and literal reconstruction of every vertical
letter; it is implemented in
\path{TNLean/MPS/MPDO/VerticalSpectral.lean}.

The positive-weight part of the third stage is also complete. For each grouped
sector $\alpha$, define
\begin{align}
    (E_\alpha)_{ab}
        &=\tr(M_\alpha^{(a,b)}).
    \label{eq:cpgsv17_vertical_diagonal_restriction_sector_trace_matrix}
\end{align}
The physical range projector of each grouped sector satisfies the required
loop identity, and horizontal separation supplies a nonzero compression. Its
trace is the sector weight $\mu_{\alpha,k}$ multiplied by a sector trace
$d_\alpha$. Under the normalization $\mu_{\alpha,1}>0$, MPDO positivity forces
\begin{align}
    d_\alpha
        &>0,
    \label{eq:cpgsv17_vertical_diagonal_restriction_positive_sector_trace}\\
    \mu_{\alpha,k}
        &>0.
    \label{eq:cpgsv17_vertical_diagonal_restriction_positive_sector_weight}
\end{align}
These results are contained in
\path{TNLean/MPS/MPDO/VerticalBNT.lean} and
\path{TNLean/MPS/MPDO/SectorTrace.lean}.

The pairwise Gram-conjugation identity represented by the second diagram at
lines~1914--1919 of Ref.~\cite{Cirac2016MPDO_arXiv} is proved in
\path{TNLean/MPS/MPDO/FigureEightPairwise.lean}. Its dependent transport to
the grouped sectors is proved in
\path{TNLean/MPS/MPDO/GroupedFigure8.lean}. Normal-tensor rigidity then gives
\begin{align}
    X_{\alpha,k}^\dagger X_{\alpha,k}
        &=
        \omega_{\alpha,k}
        X_{\alpha,1}^\dagger X_{\alpha,1}.
    \label{eq:cpgsv17_vertical_diagonal_restriction_sector_gram_scaling}
\end{align}
Here $\omega_{\alpha,k}>0$. Hence the relative gauge
$X_{\alpha,k}X_{\alpha,1}^{-1}$ is proportional to a unitary. The grouping
chooses $X_{\alpha,1}=I$. The positive scalar identity and unitary
normalization of the grouped gauge are proved in
\path{TNLean/MPS/MPDO/GroupedGramNormalization.lean}; the theorem docstrings
record the corresponding source passage at lines~1903--1921 of
Ref.~\cite{Cirac2016MPDO_arXiv}.

Projector closure and the absence of periodic vectors already yield an abstract
positive-length normal-block decomposition in
\path{TNLean/MPS/CanonicalForm/ProjectorClosureSpectral.lean}. The
isometry-preserving refinement is supplied by
\path{TNLean/MPS/MPDO/VerticalSpectral.lean}. The normal blocks are grouped
into pairwise gauge-phase-distinct BNT representatives in
\path{TNLean/MPS/MPDO/VerticalBNT.lean}. Every copy is represented by a
positive coefficient, a unit-modulus phase, and an invertible bond-space gauge,
while the physical isometries and literal vertical reconstruction are
preserved.

The unitary gauge normalizations are transported to the representative bond
spaces and composed with the physical sector isometries in
\path{TNLean/MPS/MPDO/NormalizedGroupedSectors.lean}. The
basis-of-normal-tensors property is established in
\path{TNLean/MPS/MPDO/VerticalBNTConstruction.lean}, using exact
reconstruction at positive lengths and a retained positive-dimensional
representative at length zero. The resulting orthogonal sector maps are
assembled into the final coisometry in
\path{TNLean/MPS/MPDO/VerticalCanonicalForm.lean}. Together with MPDO
positivity, the normalized BNT-refined horizontal form therefore gives both
the displayed vertical identity and exact reconstruction. Proposition~4.13
prints the former; its statement that the vertical tensor is itself in
canonical form implies the latter
\cite[Proposition~4.13]{Cirac2016MPDO_arXiv}.

The corresponding constructions under the literal CPSV canonical form are
summarized in \path{cpgsv17_vertical_cf_grouping.tex}. They establish
\path{MPOTensor.verticalCF_of_cpsvCanonicalForm} without diagonal restriction.

\paragraph{Verdict.}
Literal Proposition~4.13 is complete, and the stronger theorem under
\path{MPOTensor.IsHorizontalCF} remains available separately. Neither proof
uses diagonal restriction to construct the normal vertical representatives;
that route remains false. Proposition~4.13 does not settle the positive-fusion
problem of Appendix~C.4 or the renormalization fixed-point characterization of
Theorem~4.14 in Ref.~\cite{Cirac2016MPDO_arXiv}.

\bibliographystyle{alpha}
\bibliography{references}

\end{document}
